\documentclass[10pt]{extarticle}
\usepackage[a4paper,margin=1.5cm]{geometry}
\usepackage[T1]{fontenc}
\usepackage{amsmath,amssymb}
\usepackage{graphicx}
\usepackage[export]{adjustbox}
\usepackage{multicol}
\usepackage{enumitem}
\usepackage{booktabs,array,tabularx}
\setlength{\columnsep}{0.7cm}
\setlength{\parindent}{0pt}
\setlength{\parskip}{3pt}
\pagestyle{empty}
\makeatletter
% figure registry (figdefs.tex is generated by the extract step)
\newcommand{\qfigdef}[4]{\@namedef{qf@file@#1}{#2}\@namedef{qf@w@#1}{#3}\@namedef{qf@wide@#1}{#4}}
\newcommand{\qf@check}[1]{\@ifundefined{qf@file@#1}{\PackageError{score2}{Unknown figure id #1}{}}{}}
\newcommand{\qf@img}[4][]{\includegraphics[width=#2\dimexpr\@nameuse{qf@w@#4}pt\relax,max width=#3,max height=0.45\textheight,#1]{\@nameuse{qf@file@#4}}}
% Full-width questions: finalize wraps every question that contains a wide figure in \qfullbegin...\qfullend,
% so the WHOLE question (stem, figures, options) is set across the page in reading order, then the two columns
% resume. The question's head - its number, stem and first figure (\qfullsplit marks the end of the head) -
% is kept on one page, so a number is never left at a page foot with its figure overleaf; the rest may break
% across pages normally, so no big gaps are left. Without a \qfullsplit, the whole question is kept together
% if it is under 45% of a page. Never write these yourself.
\newif\ifqf@full
\newif\ifqf@boxed
\newsavebox{\qf@box}
\newcommand{\qf@open}{\global\qf@fulltrue\par\global\qf@boxedtrue
  \setbox\qf@box\vbox\bgroup\hsize\textwidth\linewidth\textwidth}
\newcommand{\qf@emit}[1]{\par\egroup\global\qf@boxedfalse
  \ifdim\dimexpr\ht\qf@box+\dp\qf@box\relax<#1\textheight\noindent\usebox{\qf@box}\par\else\unvbox\qf@box\fi}
\newcommand{\qfullsplit}{\ifqf@boxed\qf@emit{0.85}\fi}
\newcommand{\qf@close}{\ifqf@boxed\qf@emit{0.45}\fi\global\qf@fullfalse}
\newcommand{\qfullbegin}{\end{multicols}\qf@open}
\newcommand{\qfullend}{\qf@close\begin{multicols}{2}}
\newcommand{\qfullstart}{\qf@open}
\newcommand{\qfulldone}{\qf@close}
% \qfig[scale]{id}: block figure, centred, directly where it is written (size set by the tool).
% Fallback if a wide figure is not inside a full-width question: break out of the columns in place - no floats.
\newcommand{\qfig}[2][1]{\qf@check{#2}%
  \ifqf@full\qf@blk{#1}{#2}\else
  \ifnum\@nameuse{qf@wide@#2}=1 \expandafter\@firstoftwo\else\expandafter\@secondoftwo\fi
  {\end{multicols}\begin{center}\qf@img{#1}{\textwidth}{#2}\end{center}\begin{multicols}{2}}%
  {\qf@blk{#1}{#2}}\fi}
\newcommand{\qf@blk}[2]{\par\smallskip{\centering\qf@img{#1}{\linewidth}{#2}\par}\smallskip}
% \qfigi[scale]{id}: inline figure (inside option lists / tables / minipages); never breaks columns.
\newcommand{\qfigi}[2][1]{\qf@check{#2}\qf@img[valign=c]{#1}{\linewidth}{#2}}
\makeatother
% structure macros - lint reads these to check completeness and figure placement
\newcommand{\papersec}[1]{\par\bigskip{\centering\Large\bfseries PAPER-#1\par}\medskip}
\newcommand{\partsec}[2]{\par\medskip\noindent{\large\bfseries PART-#1: #2}\par\smallskip}
\newcommand{\sect}[1]{\par\smallskip\noindent\textbf{#1}\par}
\newcommand{\qn}[1]{\par\smallskip\noindent\textbf{#1.}~}
\newcommand{\paraq}[2]{\par\smallskip\noindent\textbf{Paragraph for Questions #1 and #2}\par}
\newenvironment{qoptions}{\begin{itemize}[leftmargin=1.8em,itemsep=1pt,topsep=2pt,label={}]}{\end{itemize}}
\newcommand{\opt}[2]{\item[(#1)] #2}
% figure registry (generated by score2_tools.py finalize)
\qfigdef{p1_phy_s1_q01_1}{TAAS - Test 01 (18 Jun 2023)-img/p1_phy_s1_q01_1.png}{90.8}{0}
\qfigdef{p1_phy_s1_q03_1}{TAAS - Test 01 (18 Jun 2023)-img/p1_phy_s1_q03_1.png}{74.0}{0}
\qfigdef{p1_phy_s1_q04_1}{TAAS - Test 01 (18 Jun 2023)-img/p1_phy_s1_q04_1.png}{70.1}{0}
\qfigdef{p1_phy_s2_q07_1}{TAAS - Test 01 (18 Jun 2023)-img/p1_phy_s2_q07_1.png}{87.0}{0}
\qfigdef{p1_phy_s2_q08_A}{TAAS - Test 01 (18 Jun 2023)-img/p1_phy_s2_q08_A.png}{78.7}{0}
\qfigdef{p1_phy_s2_q08_B}{TAAS - Test 01 (18 Jun 2023)-img/p1_phy_s2_q08_B.png}{74.7}{0}
\qfigdef{p1_phy_s2_q08_C}{TAAS - Test 01 (18 Jun 2023)-img/p1_phy_s2_q08_C.png}{56.6}{0}
\qfigdef{p1_phy_s2_q08_D}{TAAS - Test 01 (18 Jun 2023)-img/p1_phy_s2_q08_D.png}{87.4}{0}
\qfigdef{p1_phy_s2_q09_1}{TAAS - Test 01 (18 Jun 2023)-img/p1_phy_s2_q09_1.png}{99.7}{0}
\qfigdef{p1_phy_s2_q09_A}{TAAS - Test 01 (18 Jun 2023)-img/p1_phy_s2_q09_A.png}{65.9}{0}
\qfigdef{p1_phy_s2_q09_B}{TAAS - Test 01 (18 Jun 2023)-img/p1_phy_s2_q09_B.png}{66.1}{0}
\qfigdef{p1_phy_s2_q09_C}{TAAS - Test 01 (18 Jun 2023)-img/p1_phy_s2_q09_C.png}{65.9}{0}
\qfigdef{p1_phy_s2_q09_D}{TAAS - Test 01 (18 Jun 2023)-img/p1_phy_s2_q09_D.png}{55.6}{0}
\qfigdef{p1_phy_s2_q10_1}{TAAS - Test 01 (18 Jun 2023)-img/p1_phy_s2_q10_1.png}{136.2}{0}
\qfigdef{p1_phy_s2_q11_1}{TAAS - Test 01 (18 Jun 2023)-img/p1_phy_s2_q11_1.png}{181.8}{0}
\qfigdef{p1_phy_s3_q03_1}{TAAS - Test 01 (18 Jun 2023)-img/p1_phy_s3_q03_1.png}{86.7}{0}
\qfigdef{p1_phy_s3_q04_1}{TAAS - Test 01 (18 Jun 2023)-img/p1_phy_s3_q04_1.png}{100.6}{0}
\qfigdef{p1_phy_s3_q05_1}{TAAS - Test 01 (18 Jun 2023)-img/p1_phy_s3_q05_1.png}{140.1}{0}
\qfigdef{p1_phy_s3_q06_1}{TAAS - Test 01 (18 Jun 2023)-img/p1_phy_s3_q06_1.png}{79.1}{0}
\qfigdef{p1_chem_s1_q01_1}{TAAS - Test 01 (18 Jun 2023)-img/p1_chem_s1_q01_1.png}{34.6}{0}
\qfigdef{p1_chem_s1_q01_2}{TAAS - Test 01 (18 Jun 2023)-img/p1_chem_s1_q01_2.png}{37.0}{0}
\qfigdef{p1_chem_s1_q02_1}{TAAS - Test 01 (18 Jun 2023)-img/p1_chem_s1_q02_1.png}{90.3}{0}
\qfigdef{p1_chem_s1_q02_2}{TAAS - Test 01 (18 Jun 2023)-img/p1_chem_s1_q02_2.png}{92.1}{0}
\qfigdef{p1_chem_s1_q02_C}{TAAS - Test 01 (18 Jun 2023)-img/p1_chem_s1_q02_C.png}{60.7}{0}
\qfigdef{p1_chem_s1_q02_4}{TAAS - Test 01 (18 Jun 2023)-img/p1_chem_s1_q02_4.png}{130.8}{0}
\qfigdef{p1_chem_s1_q02_D}{TAAS - Test 01 (18 Jun 2023)-img/p1_chem_s1_q02_D.png}{90.9}{0}
\qfigdef{p1_chem_s1_q02_6}{TAAS - Test 01 (18 Jun 2023)-img/p1_chem_s1_q02_6.png}{94.2}{0}
\qfigdef{p1_chem_s1_q03_1}{TAAS - Test 01 (18 Jun 2023)-img/p1_chem_s1_q03_1.png}{224.3}{0}
\qfigdef{p1_chem_s2_q05_1}{TAAS - Test 01 (18 Jun 2023)-img/p1_chem_s2_q05_1.png}{55.6}{0}
\qfigdef{p1_chem_s2_q05_A}{TAAS - Test 01 (18 Jun 2023)-img/p1_chem_s2_q05_A.png}{42.4}{0}
\qfigdef{p1_chem_s2_q05_B}{TAAS - Test 01 (18 Jun 2023)-img/p1_chem_s2_q05_B.png}{52.9}{0}
\qfigdef{p1_chem_s2_q06_1}{TAAS - Test 01 (18 Jun 2023)-img/p1_chem_s2_q06_1.png}{159.6}{0}
\qfigdef{p1_chem_s2_q06_B}{TAAS - Test 01 (18 Jun 2023)-img/p1_chem_s2_q06_B.png}{47.7}{0}
\qfigdef{p1_chem_s2_q06_2}{TAAS - Test 01 (18 Jun 2023)-img/p1_chem_s2_q06_2.png}{23.7}{0}
\qfigdef{p1_chem_s2_q06_D}{TAAS - Test 01 (18 Jun 2023)-img/p1_chem_s2_q06_D.png}{33.3}{0}
\qfigdef{p1_chem_s2_q09_1}{TAAS - Test 01 (18 Jun 2023)-img/p1_chem_s2_q09_1.png}{82.0}{0}
\qfigdef{p1_chem_s3_q01_1}{TAAS - Test 01 (18 Jun 2023)-img/p1_chem_s3_q01_1.png}{137.9}{0}
\qfigdef{p1_chem_s3_q02_1}{TAAS - Test 01 (18 Jun 2023)-img/p1_chem_s3_q02_1.png}{24.0}{0}
\raggedcolumns
\makeatletter\@beginparpenalty=10000\makeatother
\newcolumntype{L}[1]{>{\raggedright\arraybackslash}p{#1}}
\begin{document}
\begin{center}{\Large\bfseries TAAS --- Test 01 (18 Jun 2023)}\end{center}
\begin{multicols}{2}
\raggedcolumns
\papersec{1}
\partsec{1}{PHYSICS}
\sect{SECTION-I(i)}
\qn{1} A cone with radius $R$ and height $H=1.0$ m is filled with water. At the bottom of the cone suddenly emerged a tiny hole with a cross section area $s=0.01\,\pi R^2$. Find time $T$ during which level of the fluid in the cone will drop by half. ($g=9.8$ m/s$^2$)
\qfig{p1_phy_s1_q01_1}
\begin{qoptions}
\opt{A}{6.43 sec}
\opt{B}{7.43 sec}
\opt{C}{8.43 sec}
\opt{D}{9.43 sec}
\end{qoptions}
\qn{2} The space between the plates of a parallel plate capacitor is filled completely consecutively with two dielectric layers 1 and 2 having the thickness $d_1$ and $d_2$ and dielectric constant $k_1$ and $k_2$ respectively. The area of each plate is equal to $A$. The density $\sigma'$ of the bound charges on the boundary plane between the two dielectrics if the voltage across the capacitor equals $V$ and electric field is directed from layer 1 to layer 2 is :
\begin{qoptions}
\opt{A}{$\left(\dfrac{k_1-k_2}{k_1+k_2}\right)\dfrac{\varepsilon_0V}{(d_1+d_2)}$}
\opt{B}{$\left(\dfrac{k_1-k_2}{k_1d_1+k_2d_2}\right)\varepsilon_0V$}
\opt{C}{$\left(\dfrac{k_2d_2-k_1d_1}{k_1+k_2}\right)\dfrac{\varepsilon_0V}{d_1d_2}$}
\opt{D}{$\left(\dfrac{k_1-k_2}{k_1d_2+k_2d_1}\right)\varepsilon_0V$}
\end{qoptions}
\qn{3} A particle of positive charge `$q$' and mass `$m$' enters a region of uniform magnetic field $B$ with velocity $V$ making an angle of $37^\circ$ with the boundary of region. Magnetic field is directed outside the plane as shown in figure. The width of the region is $d=\dfrac75\dfrac{mV}{qB}$. Find the distance travelled by the particle in the region of magnetic field.
\qfig{p1_phy_s1_q03_1}
\begin{qoptions}
\opt{A}{$\dfrac{2mV}{qB}$}
\opt{B}{$\dfrac{\pi mV}{qB}$}
\opt{C}{$\dfrac{mV}{2qB}$}
\opt{D}{$\dfrac{\pi mV}{2qB}$}
\end{qoptions}
\qn{4} Three identical thin uniform charged wires are fixed along three sides of a cube as shown in figure. Length of each wire is $l$ and linear charge density on each of them is $\lambda$. Electric field at the center of the cube is :
\qfig{p1_phy_s1_q04_1}
\begin{qoptions}
\opt{A}{0}
\opt{B}{$\dfrac{3\sqrt2\,\lambda}{4\pi\varepsilon_0l}$}
\opt{C}{$\dfrac{12\sqrt2\,\lambda}{4\pi\varepsilon_0l}$}
\opt{D}{$\dfrac{3\sqrt3\,\lambda}{16\pi\varepsilon_0l}$}
\end{qoptions}
\sect{SECTION-I(ii)}
\qn{5} A particle of mass $m$ and charge $q$ is projected from origin with velocity $\vec V=V_0\hat j$. The electric field in the region is given by

\hspace*{1em}$\vec E=E_0\hat i$ ; $y\geqslant0$

\hspace*{1em}$\vec E=0$ ; $y<0$

The magnetic field is uniform everywhere and given by $B=B_0\hat i$. Then which of the following options are correct.
\begin{qoptions}
\opt{A}{Particle moves along $x$-axis with uniform velocity.}
\opt{B}{Particle moves along $x$-axis with uniform acceleration.}
\opt{C}{The pitch of helical path followed by the particle is not constant.}
\opt{D}{Distance moved along $x$-axis after $n$ revolutions is $\dfrac n2(2n+1)\times\dfrac{\pi^2mE_0}{qB_0^2}$.}
\end{qoptions}
\qn{6} Which of the following options are correct for a small object placed very near to principal axis of a spherical mirror.
\begin{qoptions}
\opt{A}{A virtual object always have a real image in concave mirror.}
\opt{B}{A virtual object always have a virtual image in concave mirror.}
\opt{C}{A virtual object always have a real image in convex mirror.}
\opt{D}{A virtual object always have a virtual image in convex mirror.}
\end{qoptions}
\qn{7} A simple pendulum is released from rest with it's string horizontal. Different angular positions of bob are shown in the figure.
\qfig{p1_phy_s2_q07_1}
\begin{qoptions}
\opt{A}{Time taken by bob to cover arc $AC$ is greater than to arc $CD$.}
\opt{B}{Time taken by bob to cover arc $AC$ is equal to arc $CD$.}
\opt{C}{Time taken by bob to cover arc $AB$ is greater than to arc $BD$.}
\opt{D}{Time taken by bob to cover arc $AB$ is equal to arc $BD$.}
\end{qoptions}
\qn{8} Shown are four curves. Which amongst them can be possible trajectories traced by a projectile lunched from earth surface? (neglect air friction)
\begin{qoptions}
\opt{A}{\qfigi{p1_phy_s2_q08_A}}
\opt{B}{\qfigi{p1_phy_s2_q08_B}}
\opt{C}{\qfigi{p1_phy_s2_q08_C}}
\opt{D}{\qfigi{p1_phy_s2_q08_D}}
\end{qoptions}
\qn{9} In the adjacent figure, half of the space between the plates of a parallel plate capacitor is filled with dielectric material of constant $K$. Then which of the plots are possible if $X$ is measured from the plate $A$ :- ($V=$ potential, $E=$ electric field)
\qfig{p1_phy_s2_q09_1}
\begin{qoptions}
\opt{A}{\qfigi{p1_phy_s2_q09_A}}
\opt{B}{\qfigi{p1_phy_s2_q09_B}}
\opt{C}{\qfigi{p1_phy_s2_q09_C}}
\opt{D}{\qfigi{p1_phy_s2_q09_D}}
\end{qoptions}
\qn{10} In the given $R-C$ circuit switch $S$ is closed at $t=0$ and the variation of potential drop across the resistor ($V_R$) and across capacitor ($V_C$) is represented by given plots, then:
\qfig{p1_phy_s2_q10_1}
\begin{qoptions}
\opt{A}{$2V_0=2V_1+V_2$}
\opt{B}{$2V_0=V_1+V_2$}
\opt{C}{Curve I is for $V_C$ and II is for $V_R$.}
\opt{D}{Curve I is for $V_R$ and II is for $V_C$.}
\end{qoptions}
\qn{11} Figure shows a circuit in steady state. Which of the following is/are true?
\qfig{p1_phy_s2_q11_1}
\begin{qoptions}
\opt{A}{Power consumed by battery $\varepsilon_2$ is 30W.}
\opt{B}{Charge on capacitor $C$ is 40$\mu$C.}
\opt{C}{Current is resistor $R_2$ is 2A.}
\opt{D}{Current in resistor $R_3$ is 4A.}
\end{qoptions}
\qn{12} $I$ is the image of a point object $O$ formed by spherical mirror then which of the following statements is/are correct ?
\begin{qoptions}
\opt{A}{If $O$ and $I$ are on same side of the principal axis, then they have to be on opposite side of the mirror.}
\opt{B}{If $O$ and $I$ are on opposite sides of the principal axis, then they have to be on same side of the mirror.}
\opt{C}{If $O$ and $I$ are on opposite side of the principal axis, then they can be on opposite side of the mirror as well.}
\opt{D}{If $O$ and $I$ are on same side of the principal axis, then they have to be on same side of the mirror.}
\end{qoptions}
\sect{SECTION-II}
\qn{1} A small spherical uniform moon of radius $r$ and mass $m$ is orbiting circularly around a planet (uniform sphere) of mass $M$ while keeping the same face towards the planet. Any object on the moon's surface facing the planet experiences weightlessness if radius of the orbit of moon is $r\sqrt[3]{\dfrac{xM}m}$ where $x$ is \underline{\hspace{1.5cm}}
\qn{2} A concrete sphere has spherical cavity which is packed with sand dust. The relative density of concrete and sand dust are 2.4 and 0.3 respectively. For this sphere to float with its entire volume submerged under water, the ratio of mass of concrete to the mass of sand dust will be:
\qn{3} The opening near the bottom of vessel shown has an area $A$. A disk is held against the opening to keep the liquid from running out. It is given that net force on the disk applied by the liquid and air in this case is $F$. Now the disk is moved away from the opening by a short distance. The liquid comes out and strikes the disk inelastically. The initial force exerted by the liquid on the disk in this condition is $\dfrac{NF}4$. Value of $N$ is
\qfig{p1_phy_s3_q03_1}
\qn{4} In the figure shown, sphere $S_1$, $S_1$ and $S_3$ have radii $R$, $\dfrac R2$ \& $\dfrac R4$ $C_1$, $C_2$ \& $C_3$ are centers of the three spheres lying in a plane. Angle $\angle C_1C_2C_3$ is right angle. Sphere $S_3$ has a uniformly spread volume charge density $4\rho$. The remaining part of $S_2$ has uniform charge density of $2\rho$ and the left over part of $S_1$ has a uniform charge density of $\rho$. The electric field at a point $A$ at a distance $\dfrac R8$ from $C_3$ on the line $C_2C_3$ is $\dfrac{N\rho R}{\varepsilon_0}$. Find $N$
\qfig{p1_phy_s3_q04_1}
\qn{5} A free jet of an ideal fluid strikes a wedge kept on smooth horizontal surface as shown in fig. out of the total flow, a portion is deflected by $37^\circ$ and remainder is not deflected. The horizontal and vertical components of force needed to hold the wedge stationary are $F_H$ and $F_v$, respectively. Gravity is negligible. And the fluid speed remains constant. If the force ratio $\dfrac{F_H}{F_V}$ is of magnitude $1/n$ then $n=$ ?
\qfig{p1_phy_s3_q05_1}
\qn{6} Figure shows top view of a system of two semi-circular plates separated by 1 cm vertically. Vertical face in between $AF$ and $CD$ is sealed. Water is introduced radially, uniformly at rate $\pi$ lit/sec between two plates at section $ABC$. It fills all the space between plates, flows symmetrically and comes out from section $DEF$ uniformly. Internal radius ($R$) and external radius ($2R$) are 1 m and 2 m respectively. The momentum associated with flow of water in between plates is $10P$ N-sec. Find $P$.
\qfig{p1_phy_s3_q06_1}
\partsec{2}{CHEMISTRY}
\sect{SECTION-I(i)}
\qn{1}
\par\noindent\qfigi{p1_chem_s1_q01_1}$\xrightarrow[\Delta]{\mathrm{EtO^-}}$ (A) major
\par\noindent\qfigi{p1_chem_s1_q01_2}$\xrightarrow[\Delta]{{}^t\mathrm{BuO^-}\text{(tertiary butoxide ion)}}$ (B) major
\par Relation between (A) and (B) is
\begin{qoptions}
\opt{A}{Geometrical isomer}
\opt{B}{Positional isomer}
\opt{C}{Enantiomer}
\opt{D}{Chain isomer}
\end{qoptions}
\qn{2} Which of the following is incorrect?
\begin{qoptions}
\opt{A}{$\mathrm{CH_3{-}C\equiv C{-}CH_3}\xrightarrow[\mathrm{BaSO_4}]{\mathrm{D_2/Pd}}\xrightarrow{\mathrm{Br_2/CCl_4}}$\par\noindent\qfigi{p1_chem_s1_q02_1}}
\opt{B}{$\mathrm{CH_3{-}C\equiv C{-}CH_3}\xrightarrow{\mathrm{Li/Liq.\,ND_3}}\xrightarrow{\mathrm{Br_2/CCl_4}}$\par\noindent\qfigi{p1_chem_s1_q02_2}}
\opt{C}{\qfigi{p1_chem_s1_q02_C}$\xrightarrow{\mathrm{Br_2/CCl_4}}$\par\noindent\qfigi{p1_chem_s1_q02_4}}
\opt{D}{\qfigi{p1_chem_s1_q02_D} \qfigi{p1_chem_s1_q02_6}}
\end{qoptions}
\qn{3} \qfig{p1_chem_s1_q03_1}
Find maximum number of moles of A can be formed by 1 mole of acidic oxide in above reaction.
\begin{qoptions}
\opt{A}{6}
\opt{B}{2}
\opt{C}{4}
\opt{D}{8}
\end{qoptions}
\qn{4} A Pt$^{+2}$ cis complex having EAN value 84 and having two type of achiral monodentate ligand, when two same type of ligands are replaced by one butylenediamine (bn) then which of the following property is changed.
\begin{qoptions}
\opt{A}{Number of geometrical isomer}
\opt{B}{Number of optical isomer}
\opt{C}{Co-ordination number}
\opt{D}{EAN values}
\end{qoptions}
\sect{SECTION-I(ii)}
\qn{5} \qfigi{p1_chem_s2_q05_1}$\xrightarrow[\Delta]{\mathrm{NaOH}}$
\par Which of the following product can be formed in the given reaction?
\begin{qoptions}
\opt{A}{\qfigi{p1_chem_s2_q05_A}}
\opt{B}{\qfigi{p1_chem_s2_q05_B}}
\opt{C}{Et$_3$N}
\opt{D}{CH$_2=$CH$_2$}
\end{qoptions}
\columnbreak
\qn{6} \qfig{p1_chem_s2_q06_1}
\begin{qoptions}
\opt{A}{$D$ is CH$_3{-}$CH$=$CH$_2$}
\opt{B}{$E$ is \qfigi{p1_chem_s2_q06_B}}
\opt{C}{$F$ is \qfigi{p1_chem_s2_q06_2}}
\opt{D}{$B$ is \qfigi{p1_chem_s2_q06_D}}
\end{qoptions}
\qn{7} $\mathrm{2A(g)\rightarrow4B(g)+C(g)}$

At $t=0$ pressure of A is 2 atm, at $t=10$ sec. pressure of B is 2 atm and at $t=30$ sec. pressure of C is 0.75 atm. Select correct option/options :
\begin{qoptions}
\opt{A}{At $t=70$ sec pressure of B will be 3.5 atm.}
\opt{B}{Reaction may be an elementary reaction}
\opt{C}{At $t=70$ sec, rate of formation of B is $\dfrac1{80}$ atm/sec.}
\opt{D}{Reaction will complete in definite time.}
\end{qoptions}
\qn{8} A compound AB$_2$ crystallizes similarly as CaF$_2$. $r_{A^{+2}}=$ `$x$' $r_{B^-}=$ `$y$' where $r_{A^{+2}}$ and $r_{B^-}$ are respective radii. Select the correct statement(s) for the above crystal :
\begin{qoptions}
\opt{A}{Minimum distance between two B$^-$ ions will be $\dfrac2{\sqrt3}(x+y)$}
\opt{B}{Two B$^-$ ions may be present at a distance of $2(x+y)$}
\opt{C}{One A$^{+2}$ ion is surrounded by eight B$^-$ ions in lattice}
\opt{D}{It can show anion Frenkel defect}
\end{qoptions}
\qn{9} 1 mole of neon gas was taken through process $ABC$ on shown in figure. Correct option(s) for this process is/are :
\qfig{p1_chem_s2_q09_1}
(Given : $\ln\!\left(\dfrac32\right)=0.40$, $\ln\!\left(\dfrac43\right)=0.3$)
\begin{qoptions}
\opt{A}{$W_{AB}=-360$ cal}
\opt{B}{$W_{BC}=-400$ cal}
\opt{C}{$\Delta H_{AB}=1500$ cal}
\opt{D}{$q_{BC}=360$ cal}
\end{qoptions}
\qn{10} Which of the following is correct when an ideal gas having initial pressure $P$, volume $V$ and temperature $T$ is allowed to expand reversibly adiabatically until its volume becomes `$5.66\,V$' while its temperature falls to half ? (Given : $10^{0.7588}=5.66$, $\log2=0.301$)
\begin{qoptions}
\opt{A}{The gas is diatomic}
\opt{B}{The gas is monoatomic}
\opt{C}{The work done by the gas is $1.25\,PV$}
\opt{D}{Work is $-1.25\,PV$}
\end{qoptions}
\qn{11} Which option(s) is/are correct for the given metallurgical process.

$\mathrm{Ag}\xrightarrow[\mathrm{Air}]{\mathrm{NaCN_{aq.}}}\mathrm{(X)_{(aq.)}}$

$\mathrm{(X)_{(aq.)}+}$ Zinc metal $\to\mathrm{(Y)_{(aq.)}+Ag}\downarrow$
\begin{qoptions}
\opt{A}{Step-1 is called froth flotation}
\opt{B}{(X) is a diamagnetic, planar complex}
\opt{C}{(Y) is a paramagnetic, non planar complex}
\opt{D}{Same method can be used for gold.}
\end{qoptions}
\qn{12} In which option(s) pair of both reactions produces same metal containing product?
\begin{qoptions}
\opt{A}{Calcination of Malachite ;

Calcination of Calamine}
\opt{B}{$\mathrm{PbS+2PbO}\xrightarrow[\text{absence of air}]{\Delta}$ ;

$\mathrm{PbS+PbSO_4}\xrightarrow[\text{absence of air}]{\Delta}$}
\opt{C}{Roasting of Galena ;

Roasting of Sphalerite}
\opt{D}{Calcination of Calamine ;

Roasting of Sphalerite}
\end{qoptions}
\sect{SECTION-II}
\qn{1} $\mathrm{HO{-}CH_2{-}CH_2{-}\underset{\displaystyle\|}{\overset{O}{C}}{-}OEt}$\par\noindent\qfigi{p1_chem_s3_q01_1}

In the given reaction the value of $x$ is?
\qn{2} \qfigi{p1_chem_s3_q02_1}$\xrightarrow[\text{Monochlorination}]{\mathrm{Cl_2}/h\nu}$ (A) $\xrightarrow{\mathrm{Na/ether}}$ (B) $\xrightarrow[\text{monochlorination}]{\mathrm{Cl_2}/h\nu}$ total product $\xrightarrow[\text{distillation}]{\text{fractional}}$ $\underset{\text{number of fraction}}{m}$

what is the value of `$m$'?
\qn{3} In a sample of H$-$atom containing 1 mole H$-$atoms, electrons in 30\% of H$-$atoms are in 1$^{\text{st}}$ excited state and remaining are in 2$^{\text{nd}}$ excited state and ground state. If the total energy emitted when all electrons return to ground state is 642.24 kJ. Then find \% of electrons in 2$^{\text{nd}}$ excited state.

[Energy level in H$-$atom $=E_1=-13.6$eV, $E_2=-3.4$eV, $E_3=-1.5$eV, 1eV/species = 96 kJ/mole]

(Report your answer in sum of digits if answer is 32 mark your answer $3+2=5$)
\qn{4} The decomposition of A$_2$B$_2$ follows the rate equation $r=K[\mathrm{A_2B_2}]^2$.

$\mathrm{A_2B_2}\xrightarrow{K}\mathrm{A_2+B_2}$. The kinetic data of the reaction is given below.
\begin{center}\footnotesize\setlength{\tabcolsep}{3pt}
\begin{tabular}{|L{3.3cm}|L{1.1cm}|}\hline
Initial conc. of A$_2$B$_2$ (mole/lit)&7\\\hline
First half life of A$_2$B$_2$ in sec.&1440\\\hline
\end{tabular}\end{center}
If the initial concentration of A$_2$B$_2$ were 56 millimol/lit, Calculate the time (in sec) taken for the conversion of A$_2$B$_2$ to reach 60\% of initial value.

(Report your answer as $\dfrac t{1000}$ )
\qn{5} When NO reacts with aqueous solution of FeSO$_4$ then brown colour complex is formed then find the number of unpaired electron in non axial set of `d' orbital of metal.
\qn{6} In the froth flotation method of ore of a d-block metal ($M$) , ore is first suspended in mixture of oil \& water followed by treatment with CuSO$_4$ and compressed air then concentrated ore is obtained. Pure metal is obtained from concentrated ore after roasting, carbon reduction and distillation then which of the following substance given below containing that metal ($M$).

Calamine, Zinc blende, Sphalerite, Brass, Bronze, German silver
\partsec{3}{MATHEMATICS}
\sect{SECTION-I(i)}
\qn{1} If $y=\dfrac{ax+b}{cx+d}$, then value of $\dfrac{y_2^2}{y_1y_3}$, wherever exists, is (where $y_n$ denotes $n^{\text{th}}$ order derivative of differentiable function $y$ and $a,b,c,d$ are constants) -
\begin{qoptions}
\opt{A}{$\dfrac32$}
\opt{B}{$\dfrac23$}
\opt{C}{$\dfrac45$}
\opt{D}{$\dfrac54$}
\end{qoptions}
\qn{2} If $f(x)=\begin{vmatrix}(x-1)^2&(x-2)^3&(x+1)\\(x-1)^3&x^2&(x+1)^2\\x+2&x+1&x^5\end{vmatrix}$ then coefficient of $x$ in $f(x)$ is -
\begin{qoptions}
\opt{A}{$-14$}
\opt{B}{$-16$}
\opt{C}{0}
\opt{D}{25}
\end{qoptions}
\qn{3} \par\noindent\hfill$\displaystyle\begin{aligned}\int\Bigl(2e^{\tan^{-1}x+x^2+\ln x}&+\dfrac{(2x^3+2x+1)}{\left(x+\frac1x\right)}\\&\times e^{\tan^{-1}x+x^2+\ln x}\Bigr)dx\end{aligned}$\hfill\null\par\noindent is equal to
\begin{qoptions}
\opt{A}{$\dfrac{x^3}3e^{\tan^{-1}x+x^2+\ln x}+C$}
\opt{B}{$\dfrac{x^2}2e^{\tan^{-1}x+x^2+\ln x}+C$}
\opt{C}{$xe^{\tan^{-1}x+x^3+\ln x}+C$}
\opt{D}{$x^2e^{\tan^{-1}x+x^2}+C$ (where $C$ is constant of integration)}
\end{qoptions}
\qn{4} The function $f(x)=\alpha x^2-\beta x-4x^3+\gamma$, where $\alpha,\beta,\gamma\in R$, has local minima at $P(\log_2a,\,f(\log_2a))$ \& local maxima at $Q(\log_2a^2,\,f(\log_2a^2))$. If the graph of $f(x)$ changes concavity about the point $R\left(\dfrac34,\,f\left(\dfrac34\right)\right)$, then which of the following conic section can have eccentricity `$a$' -
\begin{qoptions}
\opt{A}{circle}
\opt{B}{parabola}
\opt{C}{ellipse}
\opt{D}{hyperbola}
\end{qoptions}
\sect{SECTION-I(ii)}
\qn{5} $ABC$ is a triangle in which $\angle C=2\angle B$. $D$ is a point on $BC$ such that $AD$ bisects $\angle BAC$ and $AC=BD$. Then which of the following is/are true ?
\begin{qoptions}
\opt{A}{$\cos B+\cos C=\dfrac{\sqrt5}2$}
\opt{B}{$\cos B-\cos C=\dfrac12$}
\opt{C}{$\dfrac{\sin A}{\sin C}=1$}
\opt{D}{$\dfrac{\sin A}{\sin C}=\dfrac{\sqrt5+1}4$}
\end{qoptions}
\qn{6} Which of the following statements is/are always true ?
\begin{qoptions}
\opt{A}{number of roots of equation $\displaystyle\sum_{r=1}^{2017}\dfrac r{x-r}=0$ is 2016}
\opt{B}{number of roots of equation $\displaystyle\sum_{r=1}^{2017}\dfrac r{x-r}=0$ is 2017}
\opt{C}{if $f(x)$ is continuous in $[0,2]$ such that $f(0)=f(2)$, then there exists atleast one $x\in[0,1]$ such that $f(x)=f(x+1)$.}
\opt{D}{if $f(x)$ is continuous in $[0,2]$ such that $f(0)=f(2)$, then there exists atleast one $x\in(0,1)$ such that $f(x)=f(x+1)$.}
\end{qoptions}
\qn{7} If $a+b+c=4$ and $a^2+b^2+c^2+3(ab+bc+ac)=21$ where $a,b,c\in R$, then which of the following is/are true ?
\begin{qoptions}
\opt{A}{minimum value of $(a+b)(b+c)(a+c)$ is 18}
\opt{B}{minimum value of $(a+b)(b+c)(a+c)$ is 24}
\opt{C}{maximum value of $(a+b)(b+c)(a+c)$ is $\dfrac{490}{27}$}
\opt{D}{maximum value of $(a+b)(b+c)(a+c)$ is $\dfrac{471}{27}$}
\end{qoptions}
\qn{8} Let $f(x)=\begin{cases}x[x]&-4\leqslant x<-1\\1-x&-1\leqslant x<1\\x^2-3x+2&1\leqslant x<2\\(x-2)[x]&2\leqslant x\leqslant4\end{cases}$, then (where $[.]$ denotes greatest integer function)
\begin{qoptions}
\opt{A}{number of points where $f(x)$ is non-differentiable is 6}
\opt{B}{sum of global maximum value and global minimum value of $f(x)$ is $\dfrac{63}4$}
\opt{C}{number of points of extrema of $f(x)$ is 3 points}
\opt{D}{number of solutions of $f(x)=4$ is 3}
\end{qoptions}
\qn{9} Let $F(n,x)=\displaystyle\sum_{r=1}^n\tan^{-1}\left(\dfrac{2x}{4+r(r+1)x^2}\right)$, where $x>0$ and $n\in N$ and $\dfrac d{dx}F(n,x)=G(n,x)$, then
\begin{qoptions}
\opt{A}{$F(8,2)=\tan^{-1}\left(\dfrac45\right)$}
\opt{B}{$\displaystyle\lim_{n\to\infty}F(n,2\sqrt3)=\dfrac\pi6$}
\opt{C}{If $G\left(n,\dfrac13\right)<0$, then $n\geq36$}
\opt{D}{$\displaystyle\lim_{n\to\infty}G(n,1)=-\dfrac25$}
\end{qoptions}
\qn{10} Let $f(x)$ be a non constant twice derivable function defined on $R$ such that $f(2+x)=f(2-x)$ and $f'\left(\dfrac12\right)=0=f'(1)$, then minimum number of roots of the equation $f''(x)=0$ in $(0,4)$, is
\begin{qoptions}
\opt{A}{4}
\opt{B}{3}
\opt{C}{2}
\opt{D}{5}
\end{qoptions}
\qn{11} Consider the following system of equations

$\alpha x+y+z=m$

$x+\alpha y+z=n$

$x+y+\alpha z=p$

then which of the following statement is/are correct -
\begin{qoptions}
\opt{A}{The given system of equations has no solution if $\alpha=-2$ and $m+n+p\neq0$}
\opt{B}{The given system of equations has no solution if $\alpha=1$ and $m\neq n$ or $n\neq p$ or $p\neq m$}
\opt{C}{The given system of equations has infinite solution if $\alpha=-2$ and $m+n+p=0$}
\opt{D}{The given system of equations has unique solution if $\alpha=1$ or $-2$}
\end{qoptions}
\qn{12} If the sides $a,b,c$ of a triangle are 3, 5 and 7 respectively then which of the following is/are true ?
\begin{qoptions}
\opt{A}{$\tan\left(\dfrac{A-B}2\right)=\dfrac{-1}{4\sqrt3}$}
\opt{B}{largest angle is $120^\circ$}
\opt{C}{length of median through vertex $A$ is $\dfrac{\sqrt{139}}4$}
\opt{D}{length of angle bisector through $C$ is $\dfrac{15}8$}
\end{qoptions}
\sect{SECTION-II}
\qn{1} If $\alpha,\beta,\gamma$ are the roots of $4x^3-16x^2+19x-6=0$, then the value of the determinant $\begin{vmatrix}\alpha-\beta-\gamma&2\alpha&2\alpha\\2\beta&\beta-\gamma-\alpha&2\beta\\2\gamma&2\gamma&\gamma-\alpha-\beta\end{vmatrix}$ is
\qn{2} A curve is represented by $x=\sec^2t$ and $y=\cot t$, where $t$ is the parameter. If the tangent at point $P$ on curve where $t=\dfrac\pi4$ meets the curve again at $Q$ then $|PQ|=\dfrac{p\sqrt q}2$, where $p,q$ are coprime. Find $p+q$
\qn{3} Let $f$ be a real valued function such that $f\left(\dfrac{2x-3}{x-2}\right)=5x-2$, $x\neq2$, then value of $\dfrac54f^{-1}(13)$ is
\qn{4} In a triangle $ABC$, medians $AD$ and $CE$ are drawn. If $AD=5$, $\angle DAC=\dfrac\pi8$ and $\angle ACE=\dfrac\pi4$ then the area of the triangle $ABC$ is equal to
\qn{5} Let $p_1$ and $p_2$ be two values of $p$ for which $f(x)=\begin{cases}\dfrac{x(\ell n(1+x)+\ell n(1-x))}{(\sec x-\cos x)};&x\in(-1,0)\\(p^2-3p-1)\sin x+x^2;&x\in[0,\infty)\end{cases}$ is differentiable at $x=0$ find $p_1^2+p_2^2$
\qn{6} Let $P(x)$ be a polynomial of degree 4 with $P(1)=1$, $P'(1)=0$, $P''(1)=2$, $P'''(1)=-12$ and $P''''(1)=24$ then value of $\dfrac{P(2)}{P(3)}$ is equal to
\end{multicols}
\end{document}
